% !TEX root = ../main.tex
\begin{ouverture}[son]{17}{Ondes mécaniques}{Partie 4 · Ondes et signaux}
{Onde progressive\enspace·\enspace Célérité et retard\enspace·\enspace Onde périodique\enspace·\enspace Période et longueur d'onde\enspace·\enspace Double périodicité}
\objectif{Décrire la propagation d'une perturbation mécanique dans un milieu ; distinguer onde transversale et longitudinale.}
\objectif{Exploiter la relation entre durée de propagation, distance parcourue et célérité, notamment pour localiser une source.}
\objectif{Distinguer périodicité spatiale et périodicité temporelle ; exploiter la relation $\lambda=v\,T$.}
\objectif{Déterminer les caractéristiques d'une onde à partir de représentations spatiales ou temporelles ; simuler une onde périodique.}
\prerequis{Signal sonore, propagation, vitesse de propagation, période et fréquence (classe de seconde).}
\prerequis{Fonction sinus et sa représentation graphique ; relation $d=v\,\Delta t$.}
\end{ouverture}
\essentiel
\essSection{1}{Onde mécanique progressive}
\begin{definition}[Onde mécanique progressive]
Une onde mécanique progressive est la propagation d'une \textbf{perturbation} dans un milieu matériel (corde, ressort, eau, air,
sol), \textbf{sans transport de matière} mais avec \textbf{transport d'énergie}. Elle est \textbf{transversale} si les points du milieu bougent
perpendiculairement à la direction de propagation (corde, houle), \textbf{longitudinale} s'ils bougent parallèlement (son, ressort
comprimé, ondes sismiques P).
\end{definition}
\schema{17-onde}
\essSection{2}{Célérité et retard}
\begin{loi}[Célérité d'une onde]
La perturbation parcourt la distance $d$ pendant la durée $\Delta t$ : sa \textbf{célérité} vaut $v=\dfrac{d}{\Delta t}$ (en \unit{m.s^{-1}}). Un point M situé à la distance
$d$ de la source reproduit le mouvement de celle-ci avec un \textbf{retard} $\tau=\dfrac{d}{v}$. La célérité dépend du milieu : environ
\qty{340}{m.s^{-1}} pour le son dans l'air, \qty{1500}{m.s^{-1}} dans l'eau, \qty{5900}{m.s^{-1}} dans l'acier. Le son met \qty{0,50}{s} à parcourir \qty{170}{m} dans l'air.
\end{loi}
\par\penalty0
\schema{17-celerite}
\begin{propriete}[Localiser une source]
Connaissant la célérité, la mesure d'un retard donne une distance : écho d'un sonar ou d'un capteur à ultrasons ($d=\dfrac{v\,\Delta t}{2}$
pour un aller-retour), orage (durée entre l'éclair et le tonnerre), séisme (retard des ondes S, plus lentes, sur les ondes P).
\end{propriete}
\begin{python}
d = [0.20, 0.40, 0.60, 0.80]                 # m : distances emetteur-recepteur
tau = [0.59e-3, 1.17e-3, 1.76e-3, 2.36e-3]   # s : retards mesures
v = [d[i] / tau[i] for i in range(4)]
print([round(x) for x in v])
print('moyenne :', round(sum(v) / 4), 'm/s')
\end{python}
\begin{console}
[339, 342, 341, 339]
moyenne : 340 m/s
\end{console}
\noindent\capacite{Capacité numérique} Un microcontrôleur mesure le retard d'une salve d'ultrasons pour plusieurs distances ; le programme calcule la
célérité de chaque mesure puis leur moyenne.
\essSection{3}{Ondes périodiques}
\begin{definition}[Période et longueur d'onde]
Une onde est \textbf{périodique} si la source reproduit son mouvement à intervalles réguliers. Elle présente une \textbf{double
périodicité} : temporelle, de période $T$ (en s), avec la fréquence $f=\dfrac{1}{T}$ (en Hz) ; spatiale, de période la \textbf{longueur d'onde} $\lambda$
(en m), distance entre deux points voisins qui vibrent de la même façon. Une onde est \textbf{sinusoïdale} si le mouvement de la
source est une fonction sinus du temps.
\end{definition}
\begin{loi}[Relation entre longueur d'onde, période et célérité]
Pendant une période, l'onde progresse d'une longueur d'onde :
\[\lambda=v\,T=\frac{v}{f}.\]
La fréquence est imposée par la source, la célérité par le milieu. Un la de \qty{440}{Hz} a une longueur d'onde de \qty{0,773}{m} dans l'air.
\end{loi}
\par\penalty0
\schema{17-periodique}
\begin{attention}
Ne pas confondre la vitesse des points du milieu (qui oscillent sur place) et la célérité de l'onde. Quand une onde change de milieu,
sa fréquence reste la même mais sa célérité, donc sa longueur d'onde, change.
\end{attention}
\begin{experience}[Mesurer une longueur d'onde]
Un émetteur d'ultrasons de \qty{40}{kHz} et deux récepteurs sont reliés à un oscilloscope. En éloignant l'un des récepteurs, on compte le
nombre de fois où les deux signaux redeviennent en phase sur une distance mesurée : on en déduit $\lambda$, puis $v=\lambda\,f$.
\end{experience}
\begin{rediger}
Préciser s'il s'agit d'une représentation spatiale (en fonction de $x$) ou temporelle (en fonction de $t$), mesurer plusieurs périodes
pour gagner en précision, convertir les unités (ms, kHz, cm) avant d'utiliser $\lambda=v\,T$.
\end{rediger}
\exercices

\rubrique{Découverte}
\begin{blocLivre}{ex-17-1}
\begin{exercice}[1]{La ola dans un stade}
Dans un stade, les spectateurs se lèvent puis se rassoient les uns après les autres : une « vague » fait le tour des tribunes.
\begin{questions}
\item Pourquoi peut-on parler d'onde ? Quel est le milieu de propagation ?
\item Un spectateur change-t-il de place ? Que transporte cette onde ?
\item Cette onde est-elle transversale ou longitudinale ? Justifier.
\item Citer un exemple d'onde mécanique longitudinale et un exemple d'onde transversale.
\end{questions}
\end{exercice}
\end{blocLivre}
\begin{blocLivre}{ex-17-2}
\begin{exercice}[1]{Le coup de tonnerre}
Pendant un orage, on entend le tonnerre \qty{6,0}{s} après avoir vu l'éclair.
\begin{donnees}
Célérité du son dans l'air : \qty{340}{m.s^{-1}} ; de la lumière : \qty{3,0e8}{m.s^{-1}}.
\end{donnees}
\begin{questions}
\item Calculer la distance de l'orage.
\item Calculer la durée de trajet de la lumière. Pourquoi peut-on la négliger ?
\item Justifier la règle : « compter les secondes et diviser par 3 donne la distance en kilomètres ».
\item Le son peut-il se propager dans le vide ? Et la lumière ?
\end{questions}
\end{exercice}
\end{blocLivre}
\begin{blocLivre}{ex-17-3}
\begin{exercice}[1]{Une perturbation le long d'une corde}
\schema{17-ex03}
\begin{questions}
\item Quelle distance la perturbation a-t-elle parcourue entre les deux dates ? En déduire sa célérité.
\item À quelle date le sommet de la perturbation passera-t-il au point M, situé en $x=\qty{4,0}{m}$ ?
\item La perturbation s'étend sur \qty{0,80}{m}. Combien de temps dure le mouvement d'un point de la corde ?
\item L'onde est-elle transversale ou longitudinale ?
\end{questions}
\end{exercice}
\end{blocLivre}
\begin{blocLivre}{ex-17-4}
\begin{exercice}[1]{Lire un oscillogramme}
\schema{17-ex04}
\begin{donnees}
Célérité du son dans l'air : \qty{340}{m.s^{-1}}.
\end{donnees}
\begin{questions}
\item Le signal est-il périodique ? Sinusoïdal ?
\item Mesurer la période, en mesurant plusieurs périodes pour plus de précision. En déduire la fréquence.
\item Calculer la longueur d'onde de ce son dans l'air.
\item Le son est-il audible ? (L'oreille humaine perçoit les sons de \qty{20}{Hz} à \qty{20}{kHz}.)
\end{questions}
\end{exercice}
\end{blocLivre}
\begin{blocLivre}{ex-17-5}
\begin{exercice}[1]{Un même son dans trois milieux}
Un diapason émet un la de \qty{440}{Hz}. Le son se propage dans l'air, dans l'eau et dans une barre d'acier.
\begin{donnees}
Célérités : air \qty{340}{m.s^{-1}} ; eau \qty{1500}{m.s^{-1}} ; acier \qty{5900}{m.s^{-1}}.
\end{donnees}
\begin{questions}
\item La fréquence du son dépend-elle du milieu ? Justifier.
\item Calculer la longueur d'onde dans chacun des trois milieux.
\item Dans quel milieu le son va-t-il le plus vite ? Proposer une explication.
\item Pourquoi entend-on parfois le train arriver en collant l'oreille contre le rail, avant de l'entendre dans l'air ?
\end{questions}
\end{exercice}
\end{blocLivre}
\rubrique{Entraînement}
\begin{blocLivre}{ex-17-6}
\begin{exercice}[2]{Simuler une onde périodique, en Python}
\begin{python}
from math import sin, pi
A, T, lam = 0.02, 0.50, 1.0          # m, s, m
t = 0.10                             # date choisie (s)
for i in range(9):
    x = 0.25 * i
    y = A * sin(2 * pi * (t / T - x / lam))
    print(x, 'm :', round(y * 100, 1), 'cm')
\end{python}
\begin{console}
0.0 m : 1.9 cm
0.25 m : -0.6 cm
0.5 m : -1.9 cm
0.75 m : 0.6 cm
1.0 m : 1.9 cm
1.25 m : -0.6 cm
1.5 m : -1.9 cm
1.75 m : 0.6 cm
2.0 m : 1.9 cm
\end{console}
\begin{questions}
\item \capacite{Capacité numérique}Que représentent \code{A}, \code{T} et \code{lam} ? Le programme donne-t-il une représentation spatiale ou temporelle ?
\item Quelle périodicité les valeurs affichées font-elles apparaître ? Justifier.
\item Calculer la célérité et la fréquence de l'onde.
\item Comment modifier le programme pour obtenir le mouvement du point $x=0$ au cours du temps ?
\end{questions}
\end{exercice}
\end{blocLivre}
\begin{blocLivre}{ex-17-7}
\begin{exercice}[2]{Un séisme}
\schema{17-ex07}
\begin{donnees}
Célérités : ondes P \qty{6,0}{km.s^{-1}} ; ondes S \qty{3,5}{km.s^{-1}}.
\end{donnees}
\begin{questions}
\item Repérer sur le sismogramme l'arrivée des ondes P et celle des ondes S. Lesquelles sont les plus rapides ?
\item Mesurer le retard des ondes S sur les ondes P.
\item Montrer que la distance $d$ de l'épicentre vérifie $\dfrac{d}{3{,}5}-\dfrac{d}{6{,}0}=12$ (en km et en s), puis calculer $d$.
\item Pourquoi faut-il au moins trois stations pour localiser l'épicentre ?
\end{questions}
\end{exercice}
\end{blocLivre}
\begin{blocLivre}{ex-17-8}
\begin{exercice}[2]{L'écholocation de la chauve-souris}
Une chauve-souris émet des ultrasons de \qty{40}{kHz} ; l'écho renvoyé par un insecte lui revient \qty{12}{ms} après l'émission.
\begin{donnees}
Célérité des ultrasons dans l'air : \qty{340}{m.s^{-1}}.
\end{donnees}
\begin{questions}
\item Ces ondes sont-elles audibles par l'être humain ?
\item Calculer la distance entre la chauve-souris et l'insecte.
\item Calculer la longueur d'onde des ultrasons émis.
\item Un objet plus petit que la longueur d'onde est mal détecté. Pourquoi la chauve-souris utilise-t-elle des ultrasons plutôt que des sons graves ?
\end{questions}
\end{exercice}
\end{blocLivre}
\begin{blocLivre}{ex-17-9}
\begin{exercice}[2]{Mesurer la célérité des ultrasons}
Un émetteur d'ultrasons de fréquence \qty{40,0}{kHz} et deux récepteurs sont placés en ligne. Les deux signaux reçus, observés à
l'oscilloscope, sont en phase. On éloigne le second récepteur : les signaux redeviennent en phase dix fois de suite sur une
distance de \qty{8,5}{cm}.
\begin{questions}
\item Que signifie « les deux signaux sont en phase » ?
\item En déduire la longueur d'onde des ultrasons.
\item Calculer la célérité des ultrasons dans l'air.
\item Pourquoi est-il préférable de compter dix longueurs d'onde plutôt qu'une seule ?
\end{questions}
\end{exercice}
\end{blocLivre}
\begin{blocLivre}{ex-17-10}
\begin{exercice}[2]{La houle}
Au large, une houle a une période de \qty{8,0}{s} et ses crêtes sont distantes de \qty{100}{m}.
\begin{questions}
\item Quelle grandeur correspond à la distance entre deux crêtes voisines ?
\item Calculer la célérité de la houle, en \unit{m.s^{-1}} puis en \unit{km/h}.
\item Un bouchon flotte à la surface. Décrire son mouvement. Est-il entraîné vers la côte ?
\item Combien de temps la houle met-elle pour parcourir \qty{50}{km} ?
\end{questions}
\end{exercice}
\end{blocLivre}
\begin{blocLivre}{ex-17-11}
\begin{exercice}[2]{Une double périodicité}
\schema{17-ex11}
\begin{questions}
\item Quelle représentation est spatiale ? Laquelle est temporelle ?
\item Déterminer la longueur d'onde et la période de l'onde.
\item En déduire sa célérité et sa fréquence.
\item Que vaut l'amplitude de l'onde ? Dépend-elle de la représentation choisie ?
\end{questions}
\end{exercice}
\end{blocLivre}
\rubrique{Objectif DS}
\begin{blocLivre}{ex-17-12}
\begin{exercice}[3]{L'échographie}
\begin{docu}{1}{Le principe}
Une sonde posée sur la peau émet de brèves salves d'ultrasons de \qty{3,5}{MHz}. Chaque fois qu'elles rencontrent la limite entre deux
tissus, une partie est réfléchie vers la sonde ; la durée de l'aller-retour donne la profondeur de la limite. On étale un gel entre
la sonde et la peau, car l'air réfléchit presque totalement les ultrasons.
\end{docu}
\begin{docu}{2}{Une mesure}
L'écho renvoyé par la paroi d'un organe revient \qty{65}{\micro s} après l'émission.
\end{docu}
\begin{donnees}
Célérité des ultrasons dans les tissus mous : \qty{1540}{m.s^{-1}}.
\end{donnees}
\begin{questions}
\item Calculer la longueur d'onde des ultrasons dans les tissus.
\item Calculer la profondeur de la paroi de l'organe.
\item \initiative Les détails plus petits que la longueur d'onde sont mal observés. Quel est l'intérêt d'augmenter la fréquence ?
\item À quoi sert le gel ?
\end{questions}
\end{exercice}
\end{blocLivre}
\begin{blocLivre}{ex-17-13}
\begin{exercice}[3]{Le tsunami}
\begin{docu}{1}{Une onde géante}
Un séisme sous-marin soulève brutalement le fond de l'océan : une onde, le tsunami, se propage dans tout l'océan. Sa période,
d'environ \qty{15}{min}, ne change pas au cours de la propagation.
\end{docu}
\begin{docu}{2}{Sa célérité}
Pour ces ondes très longues, la célérité ne dépend que de la profondeur $h$ de l'océan : $v=\sqrt{g\,h}$, avec $g=\qty{9,81}{N.kg^{-1}}$.
\end{docu}
\begin{questions}
\item Calculer la célérité du tsunami au large, où la profondeur est de \qty{4000}{m}, en \unit{m.s^{-1}} puis en \unit{km/h}.
\item Combien de temps met-il pour parcourir \qty{1000}{km} ?
\item Calculer sa longueur d'onde au large, puis près de la côte, où la profondeur n'est plus que de \qty{10}{m}.
\item \initiative Expliquer pourquoi la vague, presque invisible au large, devient gigantesque en arrivant sur la côte.
\end{questions}
\end{exercice}
\end{blocLivre}
\begin{blocLivre}{ex-17-14}
\begin{exercice}[3]{La profondeur d'un puits}
\begin{docu}{1}{L'expérience}
On lâche une pierre dans un puits et l'on entend le « plouf » \qty{2,5}{s} plus tard. La durée mesurée est la somme de la durée de chute
de la pierre, $t_1=\sqrt{\dfrac{2\,h}{g}}$, et de la durée de remontée du son, $t_2=\dfrac{h}{v}$.
\end{docu}
\begin{docu}{2}{Une résolution par dichotomie}
\begin{python}
g, v, t = 9.81, 340, 2.5
h_min, h_max = 0, 100
while h_max - h_min > 0.01:
    h = (h_min + h_max) / 2
    if (2 * h / g) ** 0.5 + h / v > t:
        h_max = h
    else:
        h_min = h
print(round(h, 1), 'm')
\end{python}
\begin{console}
28.6 m
\end{console}
\end{docu}
\begin{questions}
\item En négligeant la durée de remontée du son, calculer la profondeur du puits.
\item \capacite{Capacité numérique}Expliquer le principe du programme : que contient l'intervalle $[h_\mathrm{min} ; h_\mathrm{max}]$ à chaque étape ?
\item Comparer le résultat du programme à celui de la question 1. Quelle erreur relative commet-on en négligeant le son ?
\item \initiative Calculer la durée de remontée du son et vérifier le résultat.
\end{questions}
\end{exercice}
\end{blocLivre}
\begin{blocLivre}{ex-17-15}
\begin{exercice}[3]{Le capteur à ultrasons d'un robot}
\begin{docu}{1}{Le capteur}
Un capteur à ultrasons, piloté par un microcontrôleur, émet une salve et mesure la durée de l'aller-retour de l'écho. Le programme
du robot calcule la distance avec une célérité fixée à \qty{340}{m.s^{-1}}.
\end{docu}
\begin{docu}{2}{Influence de la température}
La célérité du son dans l'air dépend de la température : $v\approx331+0{,}6\,\theta$ (en \unit{m.s^{-1}}, avec $\theta$ en \unit{\degreeCelsius}).
\end{docu}
\begin{questions}
\item Calculer la célérité du son à \qty{20}{\degreeCelsius} et à \qty{0}{\degreeCelsius}.
\item Par une matinée glaciale (\qty{0}{\degreeCelsius}), l'écho revient au bout de \qty{5,83}{ms}. Calculer la distance réelle de l'obstacle et la distance affichée.
\item Calculer l'erreur relative commise. Est-elle gênante pour un robot ?
\item \initiative Quelle durée d'écho correspond à un obstacle situé à \qty{2,0}{m}, à \qty{20}{\degreeCelsius} ? Proposer une amélioration du capteur.
\end{questions}
\end{exercice}
\end{blocLivre}
